\noindent\textbf{Background: }Rapid platform change is reshaping career mobility in East Asian entertainment industries. Print media, television drama, variety shows, and creator-led video platforms differ in visibility, interaction, and crisis diffusion, so career transition must be understood as a joint adjustment of platform logic, public expectation, and personal brand resilience. The long public career of Nozomi Sasaki offers a suitable open case for template demonstration.

\noindent\textbf{Objective: }To build a thesis-ready framework for analyzing celebrity career transition, summarize stage characteristics and resilience mechanisms, and verify that common components such as figures, tables, long tables, and bilingual front matter can be rendered with fully public sample content.

\noindent\textbf{Methods: }Public reports, interviews, programme schedules, and social-media timestamps were reorganized into a three-layer framework of stage segmentation, event coding, and resilience scoring. The case was divided into launch, expansion, crisis, and restart stages. Platform shift, role diversification, public response, and self-authored communication were then compared across stages with descriptive statistics and structured exhibits.

\noindent\textbf{Results: }Three findings were highlighted. First, career transition was not a linear upgrade but a sequence of reorganizations around platform windows. Second, visibility accumulated in traditional media supported later movement into drama and variety formats, yet long-term stability depended on the ability to rebuild controllable narratives in new media after crisis. Third, image resilience could be separated into silence, repair, and restart phases, each with different communication strategies and platform mixes. All major layout elements required by the template were successfully represented with public demonstration text and graphics.

\noindent\textbf{Conclusion: }Cross-media career transition can be interpreted as a combined process of resource redistribution, narrative reconstruction, and platform reselection. A complete public demonstration thesis can therefore showcase the template without exposing any real clinical record, unpublished manuscript, or private research material.

\noindent\textbf{Keywords: }cross-media transition; case study; image resilience; East Asian entertainment industry; public demonstration template
